\subsection{Reproduction, Inheritance and Model Selection}\label{sec:inheritance}
\subsubsection{Decision-mediated paired birth}
Reproduction requires simultaneous \code{reproduce} choices by two living,
same-species, opposite-sex animals. Eligibility includes local distance,
unobstructed sight, a dry route, maturity, strict health threshold, sufficient
energy after the configured cost, cooldown, species population capacity and a
reachable empty birth cell. A parent participates in at most one birth per step.
Proximity alone never produces a child. Failed eligibility or contention leaves
successful-birth cost/cooldown unapplied, while ordinary costs and failure
feedback remain.

Maturity/cooldown variation is configured as a percentage window. For center
$c$ and percentage $p$, the integer deviation is
$\Delta=\lfloor cp/100+0.5\rfloor$; due intervals are sampled from the resulting
inclusive, positive-clipped integer window. The sampled due age is retained in
body/checkpoint state rather than redrawn on every observation. The sensing
interface may use optimistic earliest eligibility before the clock has been
sampled; physical resolution remains authoritative.

\subsubsection{Parameter crossover and rare mutation}
The child starts at the configured initial body energy and health. Both parents
pay energy; reproduction does not directly consume all parental health.
Compatible parental networks are mixed by uniform per-parameter crossover.
For a scalar coordinate $j$, the conceptual inheritance rule is
\begin{equation}
 \widetilde\theta_{c,j}=m_j\theta_{p,j}+(1-m_j)\theta_{q,j},
 \qquad m_j\sim\operatorname{Bernoulli}(1/2).
\end{equation}
The implementation traverses its algorithm-specific parameter tree and creates
independent storage. Rare mutation events choose a bounded sparse subset of
matrix weights and add configured small Gaussian perturbations. They do not
mutate the whole species or a parent's live network. Event probability, maximum
changed weights and noise scale are configurable; mutation is not an operation
on every birth and is disabled in frozen evaluation.

Crossover/mutation use deterministically derived offspring streams. The newborn
receives fresh optimizers, experience, temporal state and random generators;
Q targets or a GRPO reference are initialized consistently with the inherited
model. This is inheritance of learned weights, including acquired within-life
changes. It is an artificial mechanism and should not be described as faithful
biological genetic inheritance. Figure~\ref{fig:inheritance} separates the
mechanisms explicitly.

\subsubsection{Champion archives and run chaining}
Training tracks species/sex-specific records for health, kills, lifespan, food
energy, lineage and, for prey-eating species, a balanced hunting/survival score.
Selectors are configurable; a lifespan score includes the age of a still-living
animal and therefore need not be a completed lifespan. Models are captured at
the record-setting step after that step's learning, including terminal animals
before their controllers are retired.

Fresh experiments retain the best observed full rank tuple across runs by
strict comparison; exact ties keep the incumbent identity and parameters.
Known typed warm starts can establish the initial record. When chaining is
enabled, the selected matching species/sex record initializes the next world's
private controllers with fresh training/body state. A no-chain option instead
restores the original initialization each episode. Opt-out and legacy-resume
semantics are explicit. Optional early stopping counts completed runs without
improvement among chosen species/sex criteria; any monitored rank improvement
resets patience. This is a budget rule on observed records, not convergence or
held-out evaluation.

\subsubsection{Lineage selection}
Biological generation and model-archive generation are distinct. A newborn's
tracked founder set is the union of its parents' sets, and tracked birth
generation advances beyond the maximum known parent generation. An animal
contributes lifetime kills, actual food energy and age once to each distinct
founder family. Shared ancestry does not multiply credit within a family;
selection uses the candidate's strongest single family rather than summing
overlapping families.

Generation advantage requires maturity plus a personal kill for a prey-eating
animal, or positive food energy for a resource eater. Predator lineage ranks
then prioritize family kills; prey ranks prioritize family animal-time.
Documented tie-breakers include viable descendants, rates and personal outcomes.
Unqualified fallbacks have a provisional generation score. A qualification guard
can retain a proven starting model, and strict best-so-far retention additionally
compares the whole rank. Family ledgers reset in each new world even when weights
are chained; warm starts do not import a physical genealogy. Missing legacy
ancestry is marked partial rather than reconstructed speculatively.
